% ────────────────────────
%  FRENCH SUMMARY
% ────────────────────────
\frontchapter{Résumé}
\noindent
Ce rapport développe une chaîne complète de cartographie des types de cultures
au Rwanda à partir de séries temporelles Sentinel-2 et de modèles
d'apprentissage supervisé. L'étude porte sur la Saison~B~2025 et utilise le
jeu de données Harvard Dataverse, en considérant quatre classes de
monoculture : maïs, riz, haricot et pomme de terre irlandaise. Le travail
répond aux principales difficultés du suivi agricole en contexte de petites
exploitations, notamment la couverture nuageuse, les lacunes temporelles, le
déséquilibre des classes et le risque de fuite spatiale.

Après prétraitement, le jeu de données final contient 21\,865 QUADKEYs,
chacun représenté par une série temporelle Sentinel-2 de 30 observations. Des
indices de végétation, des variables statistiques et des descripteurs
phénologiques ont été extraits, puis réduits de 96 à 37 variables
sélectionnées. Cinq modèles classiques d'apprentissage automatique et trois
modèles profonds basés sur les séquences ont été évalués. Le modèle Extra
Trees optimisé a obtenu une exactitude de 0.9787 et un score macro F1 de
0.9305 sur le jeu de test, tandis que le modèle 1D CNN a atteint un score
macro F1 de 0.9324 comme référence complémentaire.

Le modèle final Extra Trees a été appliqué au district de Nyagatare à l'aide
d'un masque de terres cultivées Dynamic World et de seuils de confiance.
L'inférence a couvert 4\,088\,243 QUADKEYs agricoles, dont 92.2\% ont reçu
une classe de culture et 7.8\% ont été classés comme incertains. Les résultats
montrent que les séries temporelles Sentinel-2 et l'apprentissage automatique
peuvent soutenir la cartographie opérationnelle des cultures au Rwanda, tout
en soulignant la nécessité d'une validation spatiale indépendante, de tests
multi-saisons et d'une meilleure prise en compte de l'interculture.

\medskip
\noindent\textbf{Mots-clés :} cartographie des cultures, Sentinel-2,
apprentissage automatique, Rwanda, QUADKEY, sécurité alimentaire,
télédétection.
